\documentclass[aspectratio=169,11pt]{beamer}
\usetheme{default}
\setbeamertemplate{navigation symbols}{}
\setbeamertemplate{footline}{%
  \hfill\insertframenumber/\inserttotalframenumber\hspace*{4pt}\vspace{2pt}%
}
\usepackage{lmodern}
\usepackage[most]{tcolorbox}
\usepackage{listings}
\usepackage{tikz}
\usetikzlibrary{arrows.meta,positioning,shapes.geometric,fit,backgrounds}
\usepackage{booktabs}
\usepackage{hyperref}

\definecolor{lcgreen}{HTML}{2E7D32}
\definecolor{lcblue}{HTML}{1565C0}
\definecolor{lcorange}{HTML}{EF6C00}
\definecolor{lcred}{HTML}{C62828}
\definecolor{codebg}{HTML}{F5F5F5}
\definecolor{docgray}{gray}{0.55}

\newtcolorbox{conceptbox}[1][]{colback=yellow!20,colframe=yellow!60!black,title=#1,fonttitle=\bfseries,top=2pt,bottom=2pt,left=4pt,right=4pt}
\newtcolorbox{warningbox}[1][]{colback=red!15,colframe=red!60!black,title=#1,fonttitle=\bfseries,top=2pt,bottom=2pt,left=4pt,right=4pt}
\newtcolorbox{tipbox}[1][]{colback=green!10,colframe=green!50!black,title=#1,fonttitle=\bfseries,top=2pt,bottom=2pt,left=4pt,right=4pt}
\newtcolorbox{codebox}[1][]{colback=blue!10,colframe=blue!40!black,title=#1,fonttitle=\bfseries,top=2pt,bottom=2pt,left=4pt,right=4pt}

\lstset{
  basicstyle=\ttfamily\scriptsize,
  backgroundcolor=\color{codebg},
  breaklines=true,
  frame=single,
  rulecolor=\color{blue!40!black},
  keywordstyle=\color{lcblue}\bfseries,
  stringstyle=\color{lcgreen},
  commentstyle=\color{docgray}\itshape,
  language=Python,
  showstringspaces=false,
  tabsize=2,
  xleftmargin=2pt,
  xrightmargin=2pt,
  aboveskip=2pt,
  belowskip=2pt,
}

\newcommand{\docref}[1]{{\scriptsize\color{docgray}[Docs: #1]}}

\title{Error Handling \& Resilience}
\subtitle{Section 12: Building Robust Agents}
\author{LangGraph v1.0.x}
\date{March 2026}

\begin{document}

\frame{\titlepage}

\begin{frame}{Mental Model: Building Resilient Agents}
\setlength{\itemsep}{1pt}
\begin{itemize}
\item Agents run autonomously, often 24/7, in unpredictable environments
\item External APIs fail, networks timeout, data is malformed
\item Without resilience, a single failure crashes the entire workflow
\item Resilient agents gracefully handle failures and recover automatically
\item Think \emph{crash recovery}, \emph{retry logic}, \emph{fallback paths}
\end{itemize}
\vspace{0.15cm}
\begin{conceptbox}[Key Insight]
Resilience is not about preventing failures-it's about detecting them early and recovering fast.
\end{conceptbox}
\end{frame}

\begin{frame}{Why Resilience Matters}
\setlength{\itemsep}{1pt}
\begin{itemize}
\item \textbf{Long-running processes:} Agents may execute for hours or days
\item \textbf{Unreliable I/O:} Network calls, file operations, database queries can fail
\item \textbf{Cascading failures:} One bad node can corrupt downstream state
\item \textbf{Cost:} Retrying expensive operations without limits wastes tokens
\item \textbf{User experience:} Silent failures or infinite hangs degrade trust
\end{itemize}
\vspace{0.1cm}
\begin{warningbox}[Production Reality]
Without error handling, agents will fail in production. Resilience strategies are mandatory, not optional.
\end{warningbox}
\end{frame}

\begin{frame}{Node-Level Error Handling: Try/Except}
\setlength{\itemsep}{1pt}
\begin{itemize}
\item Wrap node logic in try/except blocks
\item Catch specific exceptions (APIError, TimeoutError)
\item Return error state instead of raising
\item Log error details for monitoring
\end{itemize}
\vspace{0.15cm}
{\small \texttt{def search\_node(state):}}
{\small \texttt{try: results = search\_api(...); return \{"results": results\}}}
{\small \texttt{except APIError as e: return \{"error": str(e)\}}}
\end{frame}

\begin{frame}{Retry Policy: Automatic Retries}
\setlength{\itemsep}{1pt}
\begin{itemize}
\item LangGraph supports \texttt{retry\_policy} on node definitions
\item Automatically retries a node on transient failures
\item Exponential backoff prevents overwhelming failing services
\item Configure: \texttt{max\_attempts}, \texttt{backoff\_factor}, \texttt{retry\_on}
\end{itemize}
\vspace{0.15cm}
\begin{conceptbox}[When to Use Retries]
Ideal for transient failures (network timeouts, rate limits, temporary service outages). Not suitable for logic errors or invalid input.
\end{conceptbox}
\end{frame}

\begin{frame}{Configuring RetryPolicy}
\setlength{\itemsep}{1pt}
\begin{itemize}
\item Import: \texttt{from langgraph.pregel import RetryPolicy}
\item Create: \texttt{RetryPolicy(max\_attempts=3, backoff\_factor=2.0)}
\item Specify: \texttt{retry\_on=[APIError, TimeoutError]}
\item Attach: \texttt{graph.add\_node(..., retry\_policy=rp)}
\end{itemize}
\vspace{0.1cm}
\begin{tipbox}[Backoff Timing]
Exponential backoff: 1s, 2s, 4s, 8s... Prevents overwhelming failing services.
\end{tipbox}
\end{frame}

\begin{frame}[fragile]{Error Routing: Directing Failures}
\setlength{\itemsep}{1pt}
\begin{itemize}
\item Use conditional edges to route on error state
\item Direct failed nodes to recovery or error-handling nodes
\item Example: if search fails, try fallback data source
\end{itemize}
\vspace{0.15cm}
\begin{codebox}[Routing logic]
\begin{lstlisting}
def route_on_error(state):
  if state.get("error"):
    return "fallback_search"
  return "process_results"
\end{lstlisting}
\end{codebox}
\vspace{0.1cm}
\begin{tipbox}[Tip]
Combine try/except in nodes with error routing in edges for comprehensive error handling.
\end{tipbox}
\end{frame}

\begin{frame}{Fallback Patterns: Graceful Degradation}
\setlength{\itemsep}{1pt}
\begin{itemize}
\item Primary approach fails -> use secondary (slower, cheaper, less accurate)
\item Examples: web search -> cached data; API call -> local model
\item Ensure fallback is always available and tested
\item Never degrade silently-log fallback use for monitoring
\end{itemize}
\vspace{0.15cm}
\begin{warningbox}[Fallback Risk]
If both primary and fallback fail, the error still propagates. Have a final error handler.
\end{warningbox}
\end{frame}

\begin{frame}{Durable Execution: Checkpoint Recovery}
\setlength{\itemsep}{1pt}
\begin{itemize}
\item Checkpointing saves state after each node execution
\item If agent crashes, resume from last checkpoint (not from start)
\item Dramatically reduces waste and improves reliability
\item LangGraph's default Executor uses checkpointing
\end{itemize}
\vspace{0.15cm}
\begin{conceptbox}[Durability Magic]
A 10-step process that crashes at step 8 resumes at step 8, not step 1. Recovery is instant.
\end{conceptbox}
\end{frame}

\begin{frame}{Checkpoint Configuration}
\setlength{\itemsep}{1pt}
\begin{itemize}
\item Import: \texttt{from langgraph.checkpoint.sqlite import SqliteSaver}
\item Create: \texttt{saver = SqliteSaver(db\_path="agent.db")}
\item Compile: \texttt{graph = builder.compile(checkpointer=saver)}
\item Invoke: \texttt{graph.invoke(data, config=\{"configurable": \{"thread\_id": ".."\}\})}
\end{itemize}
\vspace{0.1cm}
\begin{tipbox}[Checkpointer Options]
SqliteSaver for single-machine. PostgresSaver for distributed/production deployments.
\end{tipbox}
\end{frame}

\begin{frame}{Error Handling Flow Diagram}
\begin{center}
\begin{tikzpicture}[scale=0.9, every node/.style={font=\scriptsize}]
\node[draw,rounded rectangle,fill=blue!30] (start) at (0,0) {Node Execute};
\node[draw,diamond,aspect=3] (err) at (2,0) {Error?};
\node[draw,rectangle,fill=green!20] (success) at (4,1.5) {Next Node};
\node[draw,rectangle,fill=orange!20] (retry) at (2,-2) {Retry};
\node[draw,diamond,aspect=2.5] (retryok) at (4,-2) {Retries\\ OK?};
\node[draw,rectangle,fill=yellow!20] (fallback) at (6,-2) {Fallback};
\node[draw,rectangle,fill=red!20] (deadletter) at (8,0) {Dead Letter};
\draw[-{Latex}] (start) -- (err);
\draw[-{Latex}] (err) -- node[above] {no} (success);
\draw[-{Latex}] (err) -- node[left] {yes} (retry);
\draw[-{Latex}] (retry) -- (retryok);
\draw[-{Latex}] (retryok) -- node[above] {yes} (fallback);
\draw[-{Latex}] (fallback) -- node[below] {fails} (deadletter);
\draw[-{Latex}] (fallback) -- node[above] {ok} (success);
\end{tikzpicture}
\end{center}
\docref{Error Handling Flow}
\end{frame}

\begin{frame}{Timeout Handling for Long-Running Nodes}
\setlength{\itemsep}{1pt}
\begin{itemize}
\item Set node-level timeouts to prevent infinite waits
\item LangGraph: \texttt{timeout=seconds} on node execution
\item Catch \texttt{TimeoutError} and route to fallback or error handler
\end{itemize}
\vspace{0.15cm}
\begin{tipbox}[Timeout Strategy]
Always set reasonable timeouts on external API calls. Default to fast failure over slow hang.
\end{tipbox}
\end{frame}

\begin{frame}{Timeout Configuration Example}
\setlength{\itemsep}{1pt}
\begin{itemize}
\item Set: \texttt{graph.add\_node("api\_call", node\_fn, timeout=10)}
\item Combine with retry: \texttt{retry\_policy=RetryPolicy(...)}
\item Timeout fires -> treated as retriable error
\item Max retries enforced; fail gracefully if all attempts exhaust
\end{itemize}
\vspace{0.1cm}
\begin{tipbox}[Timeout Strategy]
Set conservative timeouts. Better fast failure than slow hang.
\end{tipbox}
\end{frame}

\begin{frame}{Dead Letter Queues: Capturing Unrecoverable Errors}
\setlength{\itemsep}{1pt}
\begin{itemize}
\item When all retries and fallbacks fail, send to "dead letter" queue
\item Preserve error state for manual review and debugging
\item Example: write to database, file, or monitoring system
\item Never silently drop errors
\end{itemize}
\vspace{0.15cm}
\begin{warningbox}[Critical]
Dead letter handling is your last defense. Ensure it's robust and monitored.
\end{warningbox}
\end{frame}

\begin{frame}{Dead Letter Pattern}
\setlength{\itemsep}{1pt}
\begin{itemize}
\item Capture failed state: \texttt{\{"timestamp", "agent\_id", "error", "state\_snapshot"\}}
\item Persist to database, file, or monitoring system
\item Never silently discard errors
\item Enable debugging and root cause analysis
\end{itemize}
\vspace{0.1cm}
\begin{warningbox}[Critical]
Dead letter is your last safety net. Ensure it's robust and monitored.
\end{warningbox}
\end{frame}

\begin{frame}{Monitoring \& Alerting with LangSmith}
\setlength{\itemsep}{1pt}
\begin{itemize}
\item LangSmith integrates with LangGraph for error tracking
\item Automatic tracing of node execution, inputs, outputs, errors
\item Set alerts on error rate thresholds
\item Dashboards for latency, token usage, retry frequency
\end{itemize}
\vspace{0.15cm}
\begin{tipbox}[Observability]
Instrument error handling with LangSmith. Without visibility, you're flying blind.
\end{tipbox}
\end{frame}

\begin{frame}{LangSmith Tracing}
\setlength{\itemsep}{1pt}
\begin{itemize}
\item Set env: \texttt{LANGSMITH\_API\_KEY}, \texttt{LANGSMITH\_PROJECT}
\item All nodes automatically traced (no code changes)
\item Logs: errors, latency, token usage, retry counts
\item Dashboards show error rates, fallback activations
\end{itemize}
\vspace{0.1cm}
\begin{tipbox}[Zero Friction]
Just set env vars; automatic observability kicks in.
\end{tipbox}
\end{frame}

\begin{frame}{Testing Error Scenarios}
\setlength{\itemsep}{1pt}
\begin{itemize}
\item Write unit tests that inject failures (mock API errors)
\item Test retry logic: verify correct number of retries
\item Test fallback activation: ensure fallback is invoked on primary failure
\item Test error routing: confirm errors reach error handlers
\item Test timeout behavior: verify timeout fires and is caught
\end{itemize}
\vspace{0.15cm}
\begin{tipbox}[Best Practice]
Every error path must be tested. Untested error handling is untested code.
\end{tipbox}
\end{frame}

\begin{frame}{Testing Retry Behavior}
\setlength{\itemsep}{1pt}
\begin{itemize}
\item Mock API to fail twice, succeed on third attempt
\item Verify result contains expected data
\item Assert call count = 3 (correct number of retries)
\item Test fallback is invoked if all retries exhaust
\end{itemize}
\vspace{0.1cm}
\begin{tipbox}[Test Pattern]
Mock external dependencies. Verify retry count, fallback activation.
\end{tipbox}
\end{frame}

\begin{frame}{Common Pitfalls}
\setlength{\itemsep}{1pt}
\begin{itemize}
\item \textbf{Swallowing errors:} Catching all exceptions, returning safe default, losing context
\item \textbf{Infinite retries:} No max\_attempts; agent retries forever on permanent failures
\item \textbf{Fallback loops:} Fallback fails too, but no error routing; state corrupted
\item \textbf{No monitoring:} Errors happen silently; no visibility into failure rates
\item \textbf{Timeout too long:} Waits 5+ minutes on failed call before failing
\end{itemize}
\vspace{0.1cm}
\docref{Error Handling Best Practices}
\end{frame}

\begin{frame}{Summary: Error Handling Strategy}
\setlength{\itemsep}{1pt}
\begin{itemize}
\item \textbf{Layer 1 (Node-level):} Try/except inside node functions
\item \textbf{Layer 2 (Automatic):} Retry policies with exponential backoff
\item \textbf{Layer 3 (Routing):} Conditional edges to fallback nodes
\item \textbf{Layer 4 (Durability):} Checkpoint recovery from crashes
\item \textbf{Layer 5 (Capture):} Dead letter handlers for unrecoverable errors
\item \textbf{Layer 6 (Visibility):} LangSmith monitoring and alerting
\end{itemize}
\vspace{0.15cm}
\begin{conceptbox}[Remember]
Resilient agents are not built; they are composed from layers of error handling.
\end{conceptbox}
\end{frame}

\begin{frame}{Cheat Sheet: Error Handling}
\small
\begin{tabular}{ll}
\toprule
\textbf{Scenario} & \textbf{Solution} \\
\midrule
Transient API failure & RetryPolicy + backoff \\
Invalid input & try/except, return error state \\
Slow external call & Set timeout, route on timeout \\
Primary source unavailable & Fallback node + conditional edge \\
Unrecoverable error & Dead letter queue + log \\
Need visibility & LangSmith tracing + alerts \\
Crash recovery & Checkpoint + thread\_id \\
\bottomrule
\end{tabular}
\end{frame}

\begin{frame}{Self-Check Questions}
\setlength{\itemsep}{2pt}
\begin{enumerate}
\item What's the difference between error routing and retry policies?
\item When would you use a fallback vs. returning an error state?
\item How does checkpointing help with resilience?
\item Why should dead letter errors be logged persistently?
\item How would you test a retry policy with exponential backoff?
\item What's the risk of setting timeouts too long?
\item How does LangSmith visibility improve error handling?
\end{enumerate}
\end{frame}

\end{document}
